%%%%%%%%%%%%%%%%%%%%%%%%%%%%%%%%%%%%%%%%%%%%%%%%%%%%%%%%%%%%%%%%%%%%%%%%%%%%%%%%
%%%%%%%%%%%%%%%%%%%%%%%%%%%%%%%%%%%%%%%%%%%%%%%%%%%%%%%%%%%%%%%%%%%%%%%%%%%%%%%%
\chapter{General Yang--Mills Gauge Theory}
\label{ch:yangmills}
%%%%%%%%%%%%%%%%%%%%%%%%%%%%%%%%%%%%%%%%%%%%%%%%%%%%%%%%%%%%%%%%%%%%%%%%%%%%%%%%
%%%%%%%%%%%%%%%%%%%%%%%%%%%%%%%%%%%%%%%%%%%%%%%%%%%%%%%%%%%%%%%%%%%%%%%%%%%%%%%%
\epigraph{%
"The laws of physics should remain invariant under local symmetry
transformations."
}{The Gauge Principle}
\section{Introduction}
The gauge principle occupies a central position in modern theoretical
physics. Electromagnetism, the weak interaction, and quantum
chromodynamics are all based upon the same underlying idea: physical
laws should be invariant under local transformations acting in an
internal space.
Historically, gauge theory developed from Maxwell's electrodynamics,
reached its modern formulation through the work of Yang and Mills, and
eventually became the mathematical foundation of the Standard Model of
particle physics. Although gauge theories are most frequently presented
within the framework of Lagrangian field theory, they possess an equally
natural Hamiltonian formulation. Since the present monograph is based on
manifestly covariant Hamiltonian mechanics, we shall emphasize both
points of view.
This chapter develops the general structure of classical Yang--Mills
theory without reference to the SHP formalism. The purpose is to
establish the universal gauge concepts that will later be extended to
five-dimensional gauge theory in Chapter~\ref{ch:shpym}. By separating
the general theory from its SHP realization, we avoid unnecessary
repetition and place the SHP construction within the broader framework
of modern gauge theory.
The presentation begins with continuous symmetry groups and their Lie
algebras. We then introduce gauge connections, covariant derivatives,
curvature tensors, Yang--Mills field equations, and the dynamics of
classical particles carrying non-Abelian charges.
%%%%%%%%%%%%%%%%%%%%%%%%%%%%%%%%%%%%%%%%%%%%%%%%%%%%%%%%%%%%%%%%%%%%%%%%%%%%%%%
\section{Continuous Symmetry Groups}
Symmetry has always played a fundamental role in physics. According to
Noether's theorem, every continuous symmetry of the action generates a
conserved quantity. In gauge theories the relevant symmetries act not on
space-time itself but on an internal vector space associated with each
space-time point.
Mathematically these symmetries are described by Lie groups.
A Lie group is simultaneously
\begin{enumerate}
\item a differentiable manifold,
\item a group,
\item and a smooth realization of the group operations.
\end{enumerate}
Typical examples encountered in physics include
\[SO(3),\]
the group of spatial rotations,
\[SO(1,3),\]
the Lorentz group,
and the compact internal groups
\[U(1),\qquad SU(2),\qquad SU(3),\]
which describe the electromagnetic, weak, and strong interactions,
respectively.
Throughout this chapter we denote a general compact Lie group by
\[G.\]
%%%%%%%%%%%%%%%%%%%%%%%%%%%%%%%%%%%%%%%%%%%%%%%%%%%%%%%%%%%%%%%%%%%%%%%%%%%%%%%
\section{Lie Algebras}
The local properties of a Lie group are determined by its Lie algebra.
Near the identity element every group element may be written as
\[\boxed{U(\theta)=\exp\left(i\theta^aT^a\right),}\]
where
\[\theta^a\]
are real parameters and
\[T^a\]
are the generators of the Lie algebra.
The generators satisfy
\[\boxed{[T^a,T^b]=if^{abc}T^c,}\]
where
\[f^{abc}\]
are the structure constants of the algebra.
These constants satisfy
\[f^{abc}=-f^{bac},\]
and obey the Jacobi identity
\[f^{abe}f^{ecd}+f^{bce}f^{ead}+f^{cae}f^{ebd}=0.\]
The Jacobi identity guarantees the consistency of successive symmetry
transformations and will later ensure the consistency of the gauge
field equations.
%%%%%%%%%%%%%%%%%%%%%%%%%%%%%%%%%%%%%%%%%%%%%%%%%%%%%%%%%%%%%%%%%%%%%%%%%%%%%%%
\section{Examples}
The three gauge groups relevant for particle physics are summarized
below.
%%%%%%%%%%%%%%%%%%%%%%%%%%%%%%%%%%%%%%%%%%%%%%%%%%%%%%%%%%%%%%%%%%%%%%%%%%%%%%%
\subsection{$U(1)$}
The Abelian group
\[U(1)=\{e^{i\alpha}\},\]
contains only one generator,
\[T=1.\]
Since there is only one generator,
\[[T,T]=0,\]
all structure constants vanish,
\[f^{abc}=0.\]
Electromagnetism is therefore an Abelian gauge theory.
%%%%%%%%%%%%%%%%%%%%%%%%%%%%%%%%%%%%%%%%%%%%%%%%%%%%%%%%%%%%%%%%%%%%%%%%%%%%%%%
\subsection{$SU(2)$}
The group
\[SU(2)\]
consists of
\[2\times2\]
unitary matrices with unit determinant.
The generators are
\[T^a=\frac12\sigma^a,\]
where
\[\sigma^a\]
are the Pauli matrices,
\[\sigma^1=\begin{pmatrix}0&1\\1&0\end{pmatrix},\]
\[\sigma^2=\begin{pmatrix}0&-i\\i&0\end{pmatrix},\]
\[\sigma^3=\begin{pmatrix}1&0\\0&-1\end{pmatrix}.\]
The commutation relations are
\[[\sigma^a,\sigma^b]=2i\epsilon^{abc}\sigma^c,\]
so that
\[f^{abc}=\epsilon^{abc}.\]
%%%%%%%%%%%%%%%%%%%%%%%%%%%%%%%%%%%%%%%%%%%%%%%%%%%%%%%%%%%%%%%%%%%%%%%%%%%%%%%
\subsection{$SU(3)$}
Quantum chromodynamics is based on
\[SU(3).\]
The generators are
\[T^a=\frac12\lambda^a,\qquad a=1,\ldots,8,\]
where
\[\lambda^a\]
are the Gell-Mann matrices.
They satisfy
\[[T^a,T^b]=if^{abc}T^c.\]
Unlike electromagnetism,
\[f^{abc}\neq0,\]
so the gauge bosons interact with one another.
This self-interaction is the fundamental origin of the nonlinear
character of Yang--Mills theory.
%%%%%%%%%%%%%%%%%%%%%%%%%%%%%%%%%%%%%%%%%%%%%%%%%%%%%%%%%%%%%%%%%%%%%%%%%%%%%%%
\section{Representations}
The generators need not act on the same vector space.
Different physical fields transform according to different
representations of the same Lie algebra.
The most important representations are
\begin{itemize}
\item fundamental representation,
\item adjoint representation,
\item tensor-product representations.
\end{itemize}
Matter fields generally belong to the fundamental representation,
while gauge fields transform in the adjoint representation,
\[A_\mu=A_\mu^aT^a.\]
This distinction will become important when deriving the transformation
laws of the gauge potential.
%%%%%%%%%%%%%%%%%%%%%%%%%%%%%%%%%%%%%%%%%%%%%%%%%%%%%%%%%%%%%%%%%%%%%%%%%%%%%%%
\section{Physical Motivation}
The mathematical formalism introduced above may at first appear
abstract. Its physical meaning is straightforward.
In ordinary mechanics, symmetry transformations act on the coordinates
of space-time. In gauge theory the transformations instead rotate the
internal state of a particle while leaving its space-time position
unchanged.
For electromagnetism the internal space is one-dimensional and the
transformation consists simply of multiplication by a phase factor.
For Yang--Mills theory the internal space has several dimensions, and
the internal orientation of a particle may vary continuously from point
to point in space-time.
The gauge field provides precisely the geometric object required to
compare these internal orientations at neighboring points. In modern
language it is a connection on a principal fiber bundle; in physical
language it specifies how internal degrees of freedom are transported
through space-time.
The remainder of this chapter develops this idea quantitatively.
%%%%%%%%%%%%%%%%%%%%%%%%%%%%%%%%%%%%%%%%%%%%%%%%%%%%%%%%%%%%%%%%%%%%%%%%%%%%%%%%
\section{Global and Local Gauge Transformations}
%%%%%%%%%%%%%%%%%%%%%%%%%%%%%%%%%%%%%%%%%%%%%%%%%%%%%%%%%%%%%%%%%%%%%%%%%%%%%%%%
The defining characteristic of a gauge theory is the requirement that
the laws of physics remain invariant under local transformations acting
in an internal vector space.
To understand the necessity of gauge fields, it is useful first to
distinguish between global and local symmetry transformations.
%%%%%%%%%%%%%%%%%%%%%%%%%%%%%%%%%%%%%%%%%%%%%%%%%%%%%%%%%%%%%%%%%%%%%%%%%%%%%%%%
\subsection{Global Transformations}
Let
\[\psi(x)\]
denote a multiplet of matter fields belonging to a representation of the
Lie group \(G\).
Under a global transformation,
\[\boxed{\psi(x)\rightarrow U\psi(x),}\]
where
\[U=\exp(i\alpha^aT^a),\]
and
\[\alpha^a=\text{constant}.\]
Since
\[\partial_\mu U=0,\]
ordinary differentiation commutes with the transformation,
\[\partial_\mu(U\psi)=U\partial_\mu\psi.\]
Consequently,
the free Lagrangian
\[\mathcal L=\bar\psi(i\gamma^\mu\partial_\mu-m)\psi\]
is invariant.
Global symmetry therefore requires no additional fields.
%%%%%%%%%%%%%%%%%%%%%%%%%%%%%%%%%%%%%%%%%%%%%%%%%%%%%%%%%%%%%%%%%%%%%%%%%%%%%%%%
\subsection{Local Transformations}
Suppose now that the symmetry parameters become functions of space-time,
\[\alpha^a=\alpha^a(x).\]
The transformation is
\[\boxed{\psi(x)\rightarrow U(x)\psi(x),}\]
with
\[U(x)=\exp(i\alpha^a(x)T^a).\]
Differentiation gives
\[\partial_\mu(U\psi)=(\partial_\mu U)\psi+U\partial_\mu\psi.\]
The additional term
\[(\partial_\mu U)\psi\]
does not vanish.
Hence
\[\partial_\mu\psi\]
does not transform in the same manner as
\[\psi.\]
The free Lagrangian therefore loses its invariance.
This simple observation is the origin of gauge theory.
%%%%%%%%%%%%%%%%%%%%%%%%%%%%%%%%%%%%%%%%%%%%%%%%%%%%%%%%%%%%%%%%%%%%%%%%%%%%%%%%
\section{The Gauge Connection}
%%%%%%%%%%%%%%%%%%%%%%%%%%%%%%%%%%%%%%%%%%%%%%%%%%%%%%%%%%%%%%%%%%%%%%%%%%%%%%%%
To restore local symmetry we introduce a new field
\[A_\mu(x),\]
called the gauge connection.
For a general Lie group,
\[\boxed{A_\mu=A_\mu^aT^a.}\]
Unlike the electromagnetic potential,
the gauge field is matrix-valued.
Its components belong to the Lie algebra of the symmetry group.
%%%%%%%%%%%%%%%%%%%%%%%%%%%%%%%%%%%%%%%%%%%%%%%%%%%%%%%%%%%%%%%%%%%%%%%%%%%%%%%%
\subsection{Covariant Derivative}
The ordinary derivative is replaced by
\[\boxed{D_\mu=\partial_\mu+igA_\mu,}\]
where
\[g\]
is the coupling constant.
The defining property of the covariant derivative is
\[\boxed{D_\mu\psi\rightarrow UD_\mu\psi.}\]
This requirement uniquely determines the transformation law of the
gauge field.
%%%%%%%%%%%%%%%%%%%%%%%%%%%%%%%%%%%%%%%%%%%%%%%%%%%%%%%%%%%%%%%%%%%%%%%%%%%%%%%%
\subsection{Transformation of the Gauge Field}
Demanding
\[D'_\mu\psi'=UD_\mu\psi\]
gives
\[(\partial_\mu+igA'_\mu)U=U(\partial_\mu+igA_\mu).\]
After straightforward algebra,
one obtains
\[\boxed{A'_\mu=UA_\mu U^{-1}-\frac{i}{g}(\partial_\mu U)U^{-1}.}\]
This is one of the fundamental equations of Yang--Mills theory.
%%%%%%%%%%%%%%%%%%%%%%%%%%%%%%%%%%%%%%%%%%%%%%%%%%%%%%%%%%%%%%%%%%%%%%%%%%%%%%%%
\subsection{Abelian Limit}
For
\[G=U(1),\]
all matrices commute,
\[U=e^{i\alpha},\]
and
\[A_\mu\]
is simply a scalar function.
The transformation law reduces immediately to
\[\boxed{A_\mu'=A_\mu-\partial_\mu\alpha,}\]
which is precisely the familiar gauge transformation of classical
electrodynamics.
Thus Maxwell theory appears naturally as the Abelian limit of the
general Yang--Mills construction.
%%%%%%%%%%%%%%%%%%%%%%%%%%%%%%%%%%%%%%%%%%%%%%%%%%%%%%%%%%%%%%%%%%%%%%%%%%%%%%%%
\section{Curvature and the Yang--Mills Field Strength}
%%%%%%%%%%%%%%%%%%%%%%%%%%%%%%%%%%%%%%%%%%%%%%%%%%%%%%%%%%%%%%%%%%%%%%%%%%%%%%%%
The introduction of the covariant derivative restores local gauge
invariance. However, it immediately raises an important question.
Is it possible to construct quantities involving the gauge field that
are independent of the particular gauge chosen?
The answer is affirmative. Such quantities arise naturally from the
noncommutativity of covariant differentiation.
%%%%%%%%%%%%%%%%%%%%%%%%%%%%%%%%%%%%%%%%%%%%%%%%%%%%%%%%%%%%%%%%%%%%%%%%%%%%%%%%
\subsection{Commutator of Covariant Derivatives}
The covariant derivative is
\[D_\mu=\partial_\mu+igA_\mu,\]
where
\[A_\mu=A_\mu^aT^a.\]
Consider two successive covariant derivatives,
\[D_\mu D_\nu\psi.\]
Using the product rule,
\[\begin{aligned}
D_\mu D_\nu\psi&=(\partial_\mu+igA_\mu)(\partial_\nu+igA_\nu)\psi\\
&=\partial_\mu\partial_\nu\psi+ig(\partial_\mu A_\nu)\psi+igA_\nu\partial_\mu\psi\\
&\quad+igA_\mu\partial_\nu\psi-g^2A_\mu A_\nu\psi.
\end{aligned}\]
Interchanging the indices gives
\[\begin{aligned}
D_\nu D_\mu\psi&=\partial_\nu\partial_\mu\psi
    +ig(\partial_\nu A_\mu)\psi+igA_\mu\partial_\nu\psi\\
&\quad+igA_\nu\partial_\mu\psi-g^2A_\nu A_\mu\psi.
\end{aligned}\]
Subtracting,
\[[D_\mu,D_\nu]\psi=D_\mu D_\nu\psi-D_\nu D_\mu\psi.\]
Since ordinary derivatives commute,
\[[\partial_\mu,\partial_\nu]=0,\]
all second derivatives cancel.
Likewise,
the mixed derivative terms cancel identically.
The remaining terms are
\[[D_\mu,D_\nu]\psi=ig(\partial_\mu A_\nu-\partial_\nu A_\mu)\psi
-g^2[A_\mu,A_\nu]\psi.\]
Extracting the common factor,
\[\boxed{[D_\mu,D_\nu]=igF_{\mu\nu},}\]
where
\[\boxed{F_{\mu\nu}=\partial_\mu A_\nu-\partial_\nu A_\mu-ig[A_\mu,A_\nu].}\]
This tensor is called the Yang--Mills field strength or curvature.
%%%%%%%%%%%%%%%%%%%%%%%%%%%%%%%%%%%%%%%%%%%%%%%%%%%%%%%%%%%%%%%%%%%%%%%%%%%%%%%%
\subsection{Component Form}
Since
\[A_\mu=A_\mu^aT^a,\]
the commutator becomes
\[[A_\mu,A_\nu]=A_\mu^aA_\nu^b[T^a,T^b].\]
Using
\[[T^a,T^b]=if^{abc}T^c,\]
one finds
\[[A_\mu,A_\nu]=if^{abc}A_\mu^aA_\nu^b T^c.\]
Consequently,
\[\boxed{F_{\mu\nu}^a=\partial_\mu A_\nu^a-
\partial_\nu A_\mu^a+gf^{abc}A_\mu^bA_\nu^c.}\]
The final nonlinear term is absent in electrodynamics and represents the
self-interaction of the gauge field.
%%%%%%%%%%%%%%%%%%%%%%%%%%%%%%%%%%%%%%%%%%%%%%%%%%%%%%%%%%%%%%%%%%%%%%%%%%%%%%%%
\subsection{Abelian Limit}
For
\[U(1),\]
all generators commute,
\[[T,T]=0,\]
and therefore
\[[A_\mu,A_\nu]=0.\]
The field tensor reduces to
\[\boxed{F_{\mu\nu}=\partial_\mu A_\nu-\partial_\nu A_\mu,}\]
which is precisely the Maxwell field tensor.
Thus classical electrodynamics appears as the Abelian limit of
Yang--Mills theory.
%%%%%%%%%%%%%%%%%%%%%%%%%%%%%%%%%%%%%%%%%%%%%%%%%%%%%%%%%%%%%%%%%%%%%%%%%%%%%%%%
\subsection{Gauge Covariance}
Under a gauge transformation,
\[A_\mu\rightarrow UA_\mu U^{-1}-\frac{i}{g}(\partial_\mu U)U^{-1},\]
one finds after direct calculation
\[\boxed{F_{\mu\nu}\rightarrow UF_{\mu\nu}U^{-1}.}\]
Unlike the gauge potential,
the field strength transforms homogeneously.
This property guarantees that traces such as
\[\operatorname{Tr}(F_{\mu\nu}F^{\mu\nu})\]
are gauge invariant and may therefore be used to construct the action.
%%%%%%%%%%%%%%%%%%%%%%%%%%%%%%%%%%%%%%%%%%%%%%%%%%%%%%%%%%%%%%%%%%%%%%%%%%%%%%%%
\subsection{Physical Interpretation}
The quantity
\[F_{\mu\nu}\]
plays exactly the same role in gauge theory as the Riemann tensor does in
general relativity.
The gauge potential
\[A_\mu\]
defines how internal states are parallel transported from one point of
space-time to another.
The curvature tensor measures the failure of this transport to be path
independent.
If
\[F_{\mu\nu}=0,\]
parallel transport is locally integrable, and the connection is locally
pure gauge.
If
\[F_{\mu\nu}\neq0,\]
parallel transport around an infinitesimal closed loop changes the
internal state of the particle.
Thus the field strength measures the intrinsic curvature of the gauge
connection.
%%%%%%%%%%%%%%%%%%%%%%%%%%%%%%%%%%%%%%%%%%%%%%%%%%%%%%%%%%%%%%%%%%%%%%%%%%%%%%%%
\section{The Yang--Mills Action}
%%%%%%%%%%%%%%%%%%%%%%%%%%%%%%%%%%%%%%%%%%%%%%%%%%%%%%%%%%%%%%%%%%%%%%%%%%%%%%%%
Having identified the gauge-invariant field strength
\[F_{\mu\nu},\]
we now seek a dynamical principle determining the evolution of the gauge
field.
As in all classical field theories, the equations of motion are obtained
from Hamilton's principle,
\[\delta S=0.\]
The principal requirement is gauge invariance.
%%%%%%%%%%%%%%%%%%%%%%%%%%%%%%%%%%%%%%%%%%%%%%%%%%%%%%%%%%%%%%%%%%%%%%%%%%%%%%%%
\subsection{Construction of the Action}
Since
\[F_{\mu\nu}\rightarrow UF_{\mu\nu}U^{-1},\]
the quantity
\[F_{\mu\nu}F^{\mu\nu}\]
is not itself gauge invariant because it is matrix-valued.
However,
the trace satisfies
\[\operatorname{Tr}(UFU^{-1})=\operatorname{Tr}(F),\]
owing to the cyclic property of the trace.
Therefore
\[\boxed{\mathcal L_{YM}=-\frac12\operatorname{Tr}(F_{\mu\nu}F^{\mu\nu})}\]
is gauge invariant.
Using the normalization
\[\operatorname{Tr}(T^aT^b)=\frac12\delta^{ab},\]
one finds
\[\boxed{\mathcal L_{YM}=-\frac14F_{\mu\nu}^aF^{a\mu\nu}.}\]
This is the Yang--Mills Lagrangian.
%%%%%%%%%%%%%%%%%%%%%%%%%%%%%%%%%%%%%%%%%%%%%%%%%%%%%%%%%%%%%%%%%%%%%%%%%%%%%%%%
\subsection{Comparison with Electrodynamics}
For
\[G=U(1),\]
the commutator term vanishes,
\[F_{\mu\nu}=\partial_\mu A_\nu-\partial_\nu A_\mu,\]
and therefore
\[\boxed{\mathcal L=-\frac14F_{\mu\nu}F^{\mu\nu},}\]
which is precisely the Maxwell Lagrangian.
Thus Maxwell theory is recovered as the Abelian limit of Yang--Mills
theory.
%%%%%%%%%%%%%%%%%%%%%%%%%%%%%%%%%%%%%%%%%%%%%%%%%%%%%%%%%%%%%%%%%%%%%%%%%%%%%%%%
\subsection{Matter Coupling}
Suppose the theory contains matter fields
\[\psi.\]
Gauge invariance requires replacing
\[\partial_\mu\rightarrow D_\mu.\]
For a Dirac field,
\[\boxed{\mathcal L_m=\bar\psi(i\gamma^\mu D_\mu-m)\psi.}\]
The complete Lagrangian becomes
\[\boxed{\mathcal L=-\frac14 F_{\mu\nu}^aF^{a\mu\nu}
+\bar\psi(i\gamma^\mu D_\mu-m)\psi.}\]
This single expression describes the interaction of matter with an
arbitrary Yang--Mills gauge field.
%%%%%%%%%%%%%%%%%%%%%%%%%%%%%%%%%%%%%%%%%%%%%%%%%%%%%%%%%%%%%%%%%%%%%%%%%%%%%%%%
\section{Variation of the Action}
We now derive the Yang--Mills equations directly from the variational
principle.
The field strength depends upon the gauge potential,
\[F_{\mu\nu}=\partial_\mu A_\nu-\partial_\nu A_\mu-ig[A_\mu,A_\nu].\]
Therefore,
the variation of the curvature is
\[\delta F_{\mu\nu}=D_\mu(\delta A_\nu)-D_\nu(\delta A_\mu),\]
where
\[D_\mu=\partial_\mu+ig[A_\mu,\cdot].\]
This remarkably compact formula replaces several pages of ordinary
component calculations.
%%%%%%%%%%%%%%%%%%%%%%%%%%%%%%%%%%%%%%%%%%%%%%%%%%%%%%%%%%%%%%%%%%%%%%%%%%%%%%%%
\subsection{Variation of the Lagrangian}
The variation is
\[\delta\mathcal L=-\frac12F^{a\mu\nu}\delta F_{\mu\nu}^a.\]
Substituting the previous result,
\[\delta\mathcal L=-F^{a\mu\nu}D_\mu(\delta A_\nu^a).\]
Integrating by parts,
and neglecting boundary terms,
one obtains
\[\delta S=\int(D_\mu F^{\mu\nu})^a\,\delta A_\nu^a\,d^4x.\]
Since
\[\delta A_\nu\]
is arbitrary,
Hamilton's principle immediately yields
\[\boxed{D_\mu F^{\mu\nu}=0.}\]
These are the vacuum Yang--Mills equations.
%%%%%%%%%%%%%%%%%%%%%%%%%%%%%%%%%%%%%%%%%%%%%%%%%%%%%%%%%%%%%%%%%%%%%%%%%%%%%%%%
\subsection{Coupling to Matter}
The matter variation contributes a current
\[J^\nu.\]
The complete equations become
\[\boxed{D_\mu F^{\mu\nu}=J^\nu.}\]
Unlike Maxwell's equations,
the derivative is covariant,
and therefore contains the gauge field itself,
\[D_\mu=\partial_\mu+ig[A_\mu,\cdot].\]
Consequently,
the Yang--Mills equations are nonlinear.
The gauge field acts as its own source.
%%%%%%%%%%%%%%%%%%%%%%%%%%%%%%%%%%%%%%%%%%%%%%%%%%%%%%%%%%%%%%%%%%%%%%%%%%%%%%%%
\subsection{Self-Interaction}
The nonlinear term
\[ig[A_\mu,F^{\mu\nu}]\]
has no analogue in classical electrodynamics.
Physically,
it implies that gauge bosons carry the very charge to which they couple.
Photons possess no electric charge,
and therefore do not interact directly.
Gluons,
however,
carry color charge,
and consequently interact with one another.
This property is responsible for the remarkable phenomena of asymptotic
freedom and color confinement in quantum chromodynamics.
%%%%%%%%%%%%%%%%%%%%%%%%%%%%%%%%%%%%%%%%%%%%%%%%%%%%%%%%%%%%%%%%%%%%%%%%%%%%%%%%
\section{Hamiltonian Formulation of Yang--Mills Theory}
%%%%%%%%%%%%%%%%%%%%%%%%%%%%%%%%%%%%%%%%%%%%%%%%%%%%%%%%%%%%%%%%%%%%%%%%%%%%%%%%
The Lagrangian formulation emphasizes manifest Lorentz covariance,
whereas the Hamiltonian formulation reveals the canonical structure of
the theory and provides the natural starting point for canonical
quantization.
Unlike ordinary particle mechanics, Yang--Mills theory is a constrained
Hamiltonian system. Not all components of the gauge potential represent
independent dynamical degrees of freedom. Some serve as Lagrange
multipliers enforcing the gauge symmetry.
The Hamiltonian formulation therefore illustrates the intimate
connection between gauge invariance and constrained dynamics.
%%%%%%%%%%%%%%%%%%%%%%%%%%%%%%%%%%%%%%%%%%%%%%%%%%%%%%%%%%%%%%%%%%%%%%%%%%%%%%%%
\subsection{Canonical Variables}
The Yang--Mills Lagrangian is
\[\mathcal L=-\frac14F_{\mu\nu}^aF^{a\mu\nu}.\]
Separating temporal and spatial indices,
\[\mu=(0,i),\qquad i=1,2,3,\]
the field tensor decomposes into
\[F_{0i},\qquad F_{ij}.\]
The canonical coordinates are
\[A_\mu^a(x).\]
Their conjugate momenta are
\[\boxed{\Pi^{a\mu}=\frac{\partial\mathcal L}{\partial(\partial_0A_\mu^a)}.}\]
%%%%%%%%%%%%%%%%%%%%%%%%%%%%%%%%%%%%%%%%%%%%%%%%%%%%%%%%%%%%%%%%%%%%%%%%%%%%%%%%
\subsection{Electric Field}
A direct calculation gives
\[\boxed{\Pi^{ai}=F^{ai0}.}\]
The canonical momentum is therefore identified with the non-Abelian
electric field,
\[E^{ai}=F^{ai0}.\]
Thus
\[\boxed{\Pi^{ai}=E^{ai}.}\]
%%%%%%%%%%%%%%%%%%%%%%%%%%%%%%%%%%%%%%%%%%%%%%%%%%%%%%%%%%%%%%%%%%%%%%%%%%%%%%%%
\subsection{Primary Constraint}
For the temporal component,
\[A_0^a,\]
the Lagrangian contains no time derivative,
\[\partial_0A_0^a.\]
Hence
\[\boxed{\Pi^{a0}=0.}\]
This is the first primary constraint.
It shows immediately that
\[A_0^a\]
is not a dynamical variable.
Instead,
it acts as a Lagrange multiplier.
%%%%%%%%%%%%%%%%%%%%%%%%%%%%%%%%%%%%%%%%%%%%%%%%%%%%%%%%%%%%%%%%%%%%%%%%%%%%%%%%
\subsection{Canonical Hamiltonian}
The Hamiltonian density is obtained by the Legendre transformation,
\[\boxed{\mathcal H=\Pi^{ai}\partial_0A_i^a-\mathcal L.}\]
After expressing
\[\partial_0A_i\]
through the electric field,
one finds
\[\boxed{\mathcal H=\frac12(E_i^aE_i^a+B_i^aB_i^a)-A_0^a(D_iE_i)^a.}\]
The first term represents the field energy,
while the second imposes Gauss's law.
%%%%%%%%%%%%%%%%%%%%%%%%%%%%%%%%%%%%%%%%%%%%%%%%%%%%%%%%%%%%%%%%%%%%%%%%%%%%%%%%
\subsection{Magnetic Field}
The magnetic field is
\[\boxed{B_i^a=\frac12\epsilon_{ijk}F_{jk}^a.}\]
Unlike electrodynamics,
\[B_i^a\]
contains quadratic terms,
\[gf^{abc}A_j^bA_k^c,\]
making the Hamiltonian nonlinear.
%%%%%%%%%%%%%%%%%%%%%%%%%%%%%%%%%%%%%%%%%%%%%%%%%%%%%%%%%%%%%%%%%%%%%%%%%%%%%%%%
\subsection{Gauss Constraint}
Variation with respect to
\[A_0^a\]
gives
\[\boxed{(D_iE_i)^a=\rho^a,}\]
where
\[\rho^a\]
is the color charge density.
In vacuum,
\[(D_iE_i)^a=0.\]
This is the non-Abelian generalization of Gauss's law.
%%%%%%%%%%%%%%%%%%%%%%%%%%%%%%%%%%%%%%%%%%%%%%%%%%%%%%%%%%%%%%%%%%%%%%%%%%%%%%%%
\subsection{Poisson Brackets}
The canonical equal-time Poisson brackets are
\[\boxed{\{A_i^a(\mathbf x),E_j^b(\mathbf y)\}=
\delta_{ij}\delta^{ab}\delta^3(\mathbf x-\mathbf y).}\]
All remaining brackets vanish.
These relations provide the classical canonical algebra of the theory.
%%%%%%%%%%%%%%%%%%%%%%%%%%%%%%%%%%%%%%%%%%%%%%%%%%%%%%%%%%%%%%%%%%%%%%%%%%%%%%%%
\subsection{Gauge Generator}
The Gauss constraint generates infinitesimal gauge transformations.
Indeed,
for
\[G[\alpha]=\int d^3x\,\alpha^a(D_iE_i)^a,\]
one finds
\[\delta A_i^a=\{A_i^a,G\},\]
which reproduces the infinitesimal Yang--Mills transformation.
Gauge symmetry therefore appears in Hamiltonian mechanics as a first-class
constraint in the sense of Dirac.
%%%%%%%%%%%%%%%%%%%%%%%%%%%%%%%%%%%%%%%%%%%%%%%%%%%%%%%%%%%%%%%%%%%%%%%%%%%%%%%%
\subsection{Physical Degrees of Freedom}
The gauge field
\[A_\mu^a\]
contains four components for each generator.
The primary constraint,
Gauss constraint,
and gauge freedom remove two components.
Consequently,
each gauge boson possesses only two propagating physical polarization
states,
exactly as in electrodynamics.
This counting remains valid before quantization and illustrates the
fundamental role of constrained Hamiltonian dynamics.
%%%%%%%%%%%%%%%%%%%%%%%%%%%%%%%%%%%%%%%%%%%%%%%%%%%%%%%%%%%%%%%%%%%%%%%%%%%%%%%%
\section{Classical Particles Carrying Non-Abelian Charge}
%%%%%%%%%%%%%%%%%%%%%%%%%%%%%%%%%%%%%%%%%%%%%%%%%%%%%%%%%%%%%%%%%%%%%%%%%%%%%%%%
Thus far the gauge field has been regarded as an independent dynamical
field. Equally important, however, is the motion of classical particles
interacting with this field.
In electrodynamics the interaction of a point particle with the
electromagnetic field is completely described by its electric charge
\(e\). Since electric charge is a scalar quantity, its value remains
constant along the particle trajectory.
The situation is fundamentally different for non-Abelian gauge theories.
The particle possesses an internal charge taking values in the Lie
algebra of the gauge group,
\[Q=Q^aT^a,\]
and this charge itself evolves during the motion.
Consequently, the complete dynamics consists of three coupled systems:
\[\begin{aligned}
    x^\mu(\tau),\\p_\mu(\tau),\\Q^a(\tau).
\end{aligned}\]
The internal degrees of freedom become dynamical variables on an equal
footing with the space-time coordinates.
%%%%%%%%%%%%%%%%%%%%%%%%%%%%%%%%%%%%%%%%%%%%%%%%%%%%%%%%%%%%%%%%%%%%%%%%%%%%%%%%
\subsection{Interaction Hamiltonian}
The gauge-covariant momentum is defined by
\[\boxed{\Pi_\mu=p_\mu-gA_\mu^aQ^a .}\]
The Hamiltonian is
\[\boxed{K=\frac{1}{2M}\Pi^\mu\Pi_\mu .}\]
Apart from replacing the electric charge by the color charge, this has
precisely the same form as the Hamiltonian for electrodynamics.
The essential difference lies in the fact that the quantities
\[Q^a\]
are themselves dynamical.
%%%%%%%%%%%%%%%%%%%%%%%%%%%%%%%%%%%%%%%%%%%%%%%%%%%%%%%%%%%%%%%%%%%%%%%%%%%%%%%%
\subsection{Hamilton Equation for the Coordinates}
Hamilton's first equation gives
\[\boxed{\dot x^\mu=\frac{\partial K}{\partial p_\mu}=\frac{\Pi^\mu}{M}.}\]
Hence
\[\Pi^\mu=M\dot x^\mu .\]
The kinetic momentum therefore determines the physical four-velocity.
%%%%%%%%%%%%%%%%%%%%%%%%%%%%%%%%%%%%%%%%%%%%%%%%%%%%%%%%%%%%%%%%%%%%%%%%%%%%%%%%
\subsection{Hamilton Equation for the Momentum}
Differentiating the Hamiltonian with respect to the coordinates,
\[\dot p_\mu=-\frac{\partial K}{\partial x^\mu},\]
one finds
\[\dot p_\mu=\frac{g}{M}\Pi^\nu Q^a\partial_\mu A_\nu^a.\]
This equation resembles the Abelian case but is not yet gauge covariant.
The covariant equation is obtained by differentiating the kinetic
momentum.
%%%%%%%%%%%%%%%%%%%%%%%%%%%%%%%%%%%%%%%%%%%%%%%%%%%%%%%%%%%%%%%%%%%%%%%%%%%%%%%%
\subsection{Evolution of the Kinetic Momentum}
Using
\[\Pi_\mu=p_\mu-gA_\mu^aQ^a,\]
one obtains
\[\frac{d\Pi_\mu}{d\tau}=\dot p_\mu-g\dot x^\nu\partial_\nu A_\mu^aQ^a
-gA_\mu^a\dot Q^a.\]
The last term shows immediately that the color charge cannot remain
constant.
A second dynamical equation governing
\[Q^a(\tau)\]
is therefore required.
%%%%%%%%%%%%%%%%%%%%%%%%%%%%%%%%%%%%%%%%%%%%%%%%%%%%%%%%%%%%%%%%%%%%%%%%%%%%%%%%
\section{Dynamics of the Color Charge}
The internal charge belongs to the Lie algebra,
\[Q=Q^aT^a.\]
Since the particle continuously changes its internal orientation while
moving through space-time, the color charge undergoes parallel transport
with respect to the gauge connection.
The natural covariant transport law is
\[\boxed{\frac{DQ}{D\tau}=0.}\]
Expanding the covariant derivative,
\[\frac{DQ}{D\tau}=\frac{dQ}{d\tau}+ig[\dot x^\mu A_\mu,Q],\]
gives
\[\boxed{\frac{dQ^a}{d\tau}=-gf^{abc}A_\mu^b\dot x^\mu Q^c.}\]
This equation expresses the precession of the color charge in internal
space.
Unlike electric charge,
color charge is not individually conserved.
Rather,
its evolution is determined by parallel transport generated by the gauge
connection.
%%%%%%%%%%%%%%%%%%%%%%%%%%%%%%%%%%%%%%%%%%%%%%%%%%%%%%%%%%%%%%%%%%%%%%%%%%%%%%%%
\section{The Wong Equations}
Substituting the color-precession equation into the evolution equation
for the kinetic momentum, one obtains after straightforward algebra
\[\boxed{\frac{d\Pi_\mu}{d\tau}=gQ^aF_{\mu\nu}^a\dot x^\nu.}\]
Since
\[\Pi^\mu=M\dot x^\mu,\]
the equation of motion becomes
\[\boxed{M\ddot x^\mu=gQ^aF^{a\mu}{}_{\nu}\dot x^\nu.}\]
Together with
\[\boxed{\frac{dQ^a}{d\tau}=-gf^{abc}A_\mu^b\dot x^\mu Q^c,}\]
these equations constitute the classical equations of motion of a
particle carrying non-Abelian charge.
They are known as the Wong equations.
%%%%%%%%%%%%%%%%%%%%%%%%%%%%%%%%%%%%%%%%%%%%%%%%%%%%%%%%%%%%%%%%%%%%%%%%%%%%%%%%
\subsection{Comparison with Electrodynamics}
For the Abelian group
\[U(1),\]
all structure constants vanish,
\[f^{abc}=0.\]
Therefore
\[\dot Q=0,\]
and the internal charge becomes a constant,
\[Q=e.\]
The Wong equations reduce immediately to
\[M\ddot x^\mu=eF^{\mu}{}_{\nu}\dot x^\nu,\]
which is precisely the covariant Lorentz-force equation.
Thus electrodynamics appears as the Abelian limit of the general
Yang--Mills dynamics.
%%%%%%%%%%%%%%%%%%%%%%%%%%%%%%%%%%%%%%%%%%%%%%%%%%%%%%%%%%%%%%%%%%%%%%%%%%%%%%%
\section{Symplectic Structure of Non-Abelian Dynamics}
%%%%%%%%%%%%%%%%%%%%%%%%%%%%%%%%%%%%%%%%%%%%%%%%%%%%%%%%%%%%%%%%%%%%%%%%%%%%%%%
The Wong equations possess a Hamiltonian formulation.
This fact is not immediately obvious because the particle evolves not
only in space-time but also in an internal color space.
Consequently,
the natural phase space is enlarged.
Instead of
\[(x^\mu,p_\mu),\]
the complete dynamical variables are
\[\boxed{(x^\mu,p_\mu,Q^a).}\]
The variables
\[Q^a\]
are not auxiliary parameters.
They are genuine dynamical coordinates describing the internal state of
the particle.
%%%%%%%%%%%%%%%%%%%%%%%%%%%%%%%%%%%%%%%%%%%%%%%%%%%%%%%%%%%%%%%%%%%%%%%%%%%%%%%
\subsection{Poisson Algebra}
The canonical variables satisfy
\[\{x^\mu,p_\nu\}=\delta^\mu_\nu .\]
The color variables satisfy
\[\boxed{\{Q^a,Q^b\}=f^{abc}Q^c.}\]
This Poisson algebra reproduces the Lie algebra of the gauge group.
Consequently,
the internal color space possesses its own symplectic structure.
%%%%%%%%%%%%%%%%%%%%%%%%%%%%%%%%%%%%%%%%%%%%%%%%%%%%%%%%%%%%%%%%%%%%%%%%%%%%%%%
\subsection{Extended Phase Space}
The total phase space is therefore
\[\boxed{\Gamma=T^*M\times\mathcal O,}\]
where
\[\mathcal O\]
is the coadjoint orbit of the Lie group.
Thus space-time motion and internal color dynamics
are unified into a single Hamiltonian system.
%%%%%%%%%%%%%%%%%%%%%%%%%%%%%%%%%%%%%%%%%%%%%%%%%%%%%%%%%%%%%%%%%%%%%%%%%%%%%%%
\subsection{Hamiltonian}
The Hamiltonian is
\[\boxed{K=\frac1{2M}(p_\mu-gA_\mu^aQ^a)^2.}\]
Using the extended Poisson brackets,
Hamilton's equations become
\[\dot z^A=\{z^A,K\},\]
where
\[z^A=(x^\mu,p_\mu,Q^a).\]
One immediately recovers
Hamilton's equations,
the Wong equations,
and the color-precession equation.
Thus the complete classical Yang--Mills dynamics is Hamiltonian.
%%%%%%%%%%%%%%%%%%%%%%%%%%%%%%%%%%%%%%%%%%%%%%%%%%%%%%%%%%%%%%%%%%%%%%%%%%%%%%%
\subsection{Geometrical Interpretation}
The phase space possesses two independent geometric structures.
The first is the ordinary symplectic structure
\[dp_\mu\wedge dx^\mu,\]
describing motion through space-time.
The second is the Kirillov--Kostant--Souriau symplectic structure on the
coadjoint orbit,
which governs the evolution of the color charge.
The interaction Hamiltonian couples these two geometries through the
gauge connection.
The resulting Hamiltonian flow simultaneously transports the particle in
space-time and rotates its internal state.
